% shader-operator-trajectory.tex
% Shader-Operator Trajectories in Neural Partition Space:
% Real-Time GPU Observation of Regime Dynamics via the
% Rendering-Measurement Identity
% Author: Kundai Farai Sachikonye
% Date: 2026

\documentclass[11pt,twocolumn,a4paper]{article}

% --- Packages ---
\usepackage[utf8]{inputenc}
\usepackage[T1]{fontenc}
\usepackage{lmodern}
\usepackage{amsmath,amssymb,amsthm,mathtools}
\usepackage{physics}
\usepackage{bm}
\usepackage{graphicx}
\usepackage{xcolor}
\usepackage{booktabs}
\usepackage{array}
\usepackage{multirow}
\usepackage{enumitem}
\usepackage{float}
\usepackage{caption}
\usepackage{subcaption}
\usepackage[colorlinks=true,linkcolor=blue!60!black,citecolor=blue!60!black,urlcolor=blue!60!black]{hyperref}
\usepackage[margin=1.8cm,top=2.2cm,bottom=2.2cm]{geometry}
\usepackage{natbib}
\usepackage{fancyhdr}
\usepackage{titlesec}
\usepackage{microtype}
\usepackage{listings}
\usepackage{tabularx}

% --- Listings Style for GLSL ---
\lstdefinestyle{glsl}{
  language=C,
  basicstyle=\ttfamily\scriptsize,
  keywordstyle=\color{blue!70!black}\bfseries,
  commentstyle=\color{green!50!black}\itshape,
  stringstyle=\color{red!60!black},
  numbers=left,
  numberstyle=\tiny\color{gray},
  numbersep=4pt,
  frame=single,
  breaklines=true,
  columns=fullflexible,
  morekeywords={vec2,vec3,vec4,mat2,mat3,mat4,sampler2D,
    uniform,varying,in,out,layout,location,precision,
    highp,mediump,lowp,texture,gl_FragColor,void,float,int,
    bool,ivec2,ivec3,ivec4,uint,uvec2,uvec3,uvec4,
    gl_Position,gl_VertexID,gl_FragCoord,discard,
    flat,smooth,centroid,const,struct,return,if,else,for}
}

% --- Theorem Environments ---
\newtheorem{theorem}{Theorem}[section]
\newtheorem{lemma}[theorem]{Lemma}
\newtheorem{corollary}[theorem]{Corollary}
\newtheorem{proposition}[theorem]{Proposition}
\newtheorem{definition}[theorem]{Definition}
\newtheorem{remark}[theorem]{Remark}
\newtheorem{axiom}{Axiom}
\newtheorem{example}[theorem]{Example}

% --- Custom Commands ---
\newcommand{\kb}{k_{\mathrm{B}}}
\newcommand{\Sk}{S_k}
\newcommand{\St}{S_t}
\newcommand{\Se}{S_e}
\newcommand{\BigO}{\mathcal{O}}
\newcommand{\R}{\mathbb{R}}
\newcommand{\N}{\mathbb{N}}
\newcommand{\Lag}{\mathcal{L}}
\newcommand{\Pot}{\Phi}
\newcommand{\qvec}{\mathbf{q}}
\newcommand{\Svec}{\mathbf{S}}
\newcommand{\Kc}{K_c}
\newcommand{\sigmin}{\sigma^2_{\mathrm{min}}}
\newcommand{\Rend}{\mathcal{R}}
\newcommand{\Meas}{\mathcal{M}}
\newcommand{\Obs}{\mathcal{O}}
\newcommand{\Tex}{\mathcal{T}}
\newcommand{\dcat}{d_{\mathrm{cat}}}
\DeclareMathOperator{\Tr}{Tr}
\DeclareMathOperator{\Vis}{Vis}
\DeclareMathOperator{\sgn}{sgn}

% --- Page Style ---
\pagestyle{fancy}
\fancyhf{}
\fancyhead[L]{\small\textit{Shader-Operator Trajectories in Neural Partition Space}}
\fancyhead[R]{\small\thepage}
\renewcommand{\headrulewidth}{0.4pt}

% --- Section Formatting ---
\titleformat{\section}{\large\bfseries}{\thesection.}{0.5em}{}
\titleformat{\subsection}{\normalsize\bfseries}{\thesubsection.}{0.5em}{}
\titleformat{\subsubsection}{\normalsize\itshape}{\thesubsubsection.}{0.5em}{}

% ============================================================================
\begin{document}

\twocolumn[
\begin{@twocolumnfalse}
\begin{center}

{\LARGE\bfseries Shader-Operator Trajectories in\\[4pt]
Neural Partition Space:\\[4pt]
Real-Time GPU Observation of Regime Dynamics\\[4pt]
via the Rendering--Measurement Identity}

\vspace{12pt}

{\large Kundai Farai Sachikonye}

\vspace{6pt}

{\small\itshape AIMe Registry for Artificial Intelligence \\[0.3em]
\texttt{kundai.sachikonye@bitspark.com}}

\vspace{6pt}

{\small April 2026}

\vspace{14pt}

\begin{minipage}{0.92\textwidth}
\textbf{Abstract.}
We derive the complete GPU implementation of real-time neural partition observation from three axioms: bounded phase space ($\mu(\Omega) < \infty$), no null state, and finite observational resolution ($\delta > 0$).
The Rendering--Measurement Identity~\cite{sachikonye2026gpu} establishes that a fragment shader evaluating partition coordinates at each texel \emph{is} a measurement, not a visualization of one; we apply this identity to the neural domain, where the Kuramoto order parameter $R$, the phase variance $\sigma^2$, and the S-entropy coordinates $(\Sk, \St, \Se) \in [0,1]^3$ constitute the five-dimensional generalized coordinate space $\qvec = (R, \sigma^2, \Sk, \St, \Se)$.
We specify the neural domain \texttt{computePartitionValue()} shader function, prove that the structural factor $S = R \exp(-\sigma^2 / 2\pi^2)$ can be evaluated as a per-texel GPU operation, and implement the Euler--Lagrange gradient flow equations~\cite{sachikonye2026lagrangian} as a single-pass integration shader.
Interference-based regime matching---superposing a query neural state against known regime signatures as wave amplitudes---replaces algorithmic comparison with physical interference, where visibility $\mathcal{V} \geq 0.95$ indicates regime identification.
We prove that integrated GPUs ($\sim\!25$~MB working set, $\sim\!1$--$2$~TFLOPS) suffice for the complete pipeline: S-entropy computation from multichannel EEG at 256~Hz, regime classification into the five operational regimes (turbulent, aperture-dominated, hierarchical cascade, coherent, phase-locked), trajectory integration, and Poincar\'{e} backward computation from specified terminus states.
Applications include real-time sleep stage detection as regime sequences, drug-induced regime traversal monitoring, and consciousness window estimation via the intersection $C(t) = P_{\mathrm{decay}} \cap T_{\mathrm{decay}}$.
All operations execute at $>30$~fps on commodity hardware.

\vspace{8pt}
\noindent\textbf{Keywords:} neural partition, fragment shader, Kuramoto order parameter, S-entropy, regime classification, GPU observation, Euler--Lagrange integration, interference matching, Poincar\'{e} computing, EEG, sleep staging
\end{minipage}

\vspace{18pt}
\end{center}
\end{@twocolumnfalse}
]

% ============================================================================
%                    SECTION 1: INTRODUCTION
% ============================================================================

\section{Introduction}
\label{sec:introduction}

Neural systems are bounded, finite, and oscillatory, yet they compute with extraordinary fidelity~\cite{buzsaki2006,fries2005}.
The neural partition framework~\cite{sachikonye2026regime,sachikonye2026aperture,sachikonye2026operator,sachikonye2026purpose} classifies brain states through five operational regimes parameterized by the Kuramoto order parameter $R$~\cite{kuramoto1975,strogatz2000,acebron2005}: turbulent ($R < 0.3$), aperture-dominated ($0.3 \leq R < 0.5$), hierarchical cascade ($0.5 \leq R < 0.8$), coherent ($0.8 \leq R < 0.95$), and phase-locked ($R \geq 0.95$).
The Euler--Lagrange equations derived from the Onsager--Machlup action on the five-dimensional coordinate space $\qvec = (R, \sigma^2, \Sk, \St, \Se)$~\cite{sachikonye2026lagrangian} govern transitions between these regimes.
The Rendering--Measurement Identity~\cite{sachikonye2026gpu} proves that a fragment shader performing categorical observation \emph{is} a measurement apparatus---the rendered texture is the measurement result, not a picture of it.

The present paper closes the loop: we specify the \emph{neural domain shader}, the concrete fragment shader implementation that observes neural partition states on the GPU.
The contribution is threefold.

\begin{enumerate}[leftmargin=*,itemsep=2pt]
\item \textbf{Neural \texttt{computePartitionValue()}.} We derive the partition observation function for the neural domain from the three axioms and implement it in GLSL~ES~3.00~\cite{khronos2023glsl}.
The shader takes S-entropy uniforms $(\Sk, \St, \Se)$ computed from multichannel EEG/LFP recordings on the CPU and evaluates the neural partition value at each texel.

\item \textbf{GPU Euler--Lagrange integration.} The gradient flow equations $\dot{R}$, $\dot{\sigma}^2$, $\dot{\Sk}$, $\dot{\St}$, $\dot{\Se}$ from the variational principle~\cite{sachikonye2026lagrangian} are integrated in a single-pass shader, enabling real-time trajectory computation through regime space.

\item \textbf{Interference-based regime matching.} By superposing a query observation texture against stored regime signature textures using the universal interference shader, regime identification reduces to reading interference visibility---a physical observable, not an algorithmic comparison.
\end{enumerate}

We prove that the entire pipeline executes within $\sim\!25$~MB GPU memory at $>30$~fps on integrated GPUs (Intel UHD~620, AMD Radeon Vega~8, Apple M1), eliminating discrete GPU, cloud, and datacenter requirements.

Section~\ref{sec:axioms} recapitulates the three axioms and derives the neural partition coordinates.
Section~\ref{sec:sentropy} specifies S-entropy computation from EEG/LFP data.
Section~\ref{sec:kuramoto} derives the Kuramoto order parameter from multichannel recordings.
Section~\ref{sec:shader} presents the neural domain shader with complete GLSL code.
Section~\ref{sec:structural} derives the structural factor computation.
Section~\ref{sec:euler-lagrange} implements the Euler--Lagrange gradient flow on GPU.
Section~\ref{sec:interference} develops interference-based regime matching.
Section~\ref{sec:observables} defines the GPU physical observables.
Section~\ref{sec:poincare} presents Poincar\'{e} computing on GPU.
Section~\ref{sec:applications} gives applications.
Section~\ref{sec:hardware} proves integrated GPU sufficiency.
Section~\ref{sec:conclusion} concludes.


% ============================================================================
%                    SECTION 2: AXIOMS AND NEURAL PARTITION COORDINATES
% ============================================================================

\section{Axioms and Neural Partition Coordinates}
\label{sec:axioms}

\begin{axiom}[Bounded Phase Space]
\label{ax:bounded}
Every persistent neural system occupies a bounded region $\Gamma \subset \R^{2f}$ of phase space with finite Liouville measure: $\mu(\Gamma) < \infty$.
\end{axiom}

\begin{axiom}[No Null State]
\label{ax:nonull}
There exists no state $x \in \Gamma$ at which all dynamical variables vanish simultaneously.
The system oscillates from categorical necessity~\cite{sachikonye2026gpu}.
\end{axiom}

\begin{axiom}[Finite Observational Resolution]
\label{ax:resolution}
Every observer has finite resolution $\delta > 0$, inducing a partition of $\Gamma$ into finitely many distinguishable cells.
\end{axiom}

From these three axioms, the neural partition coordinates $(n, \ell, m, s)$ with shell capacity $C(n) = 2n^2$ follow by hierarchical subdivision of the bounded spherical phase space~\cite{sachikonye2026operator}.
The Triple Equivalence Theorem establishes:

\begin{equation}
\boxed{S_{\mathrm{osc}} = S_{\mathrm{cat}} = S_{\mathrm{part}} = \kb M \ln n}
\label{eq:triple}
\end{equation}

\noindent where $M$ is the number of independent degrees of freedom and $n$ the number of distinguishable states per degree of freedom.

\begin{definition}[Neural Partition State]
\label{def:neural-partition}
A neural partition state is specified by the five-dimensional generalized coordinate vector
\begin{equation}
\qvec = (R, \sigma^2, \Sk, \St, \Se) \in \mathcal{M},
\label{eq:gencoords}
\end{equation}
where $\mathcal{M} = [0,1] \times [0, 2\pi^2] \times [0,1]^3$ is the compact configuration manifold.
Here $R$ is the Kuramoto order parameter, $\sigma^2$ the circular phase variance, and $(\Sk, \St, \Se) \in [0,1]^3$ the S-entropy coordinates.
\end{definition}

\begin{definition}[Operational Regimes]
\label{def:regimes}
The five operational regimes are classified by $R$:
\begin{center}
\begin{tabular}{lll}
\toprule
Regime & $R$ range & Character \\
\midrule
Turbulent & $R < 0.3$ & Desynchronized \\
Aperture-dominated & $0.3 \leq R < 0.5$ & Selective gating \\
Hierarchical cascade & $0.5 \leq R < 0.8$ & Multi-scale \\
Coherent & $0.8 \leq R < 0.95$ & Global sync \\
Phase-locked & $R \geq 0.95$ & Rigid lock \\
\bottomrule
\end{tabular}
\end{center}
The regime boundaries are determined by the bifurcation structure of the Kuramoto mean-field equations~\cite{kuramoto1975,acebron2005}.
\end{definition}


% ============================================================================
%                    SECTION 3: S-ENTROPY FROM EEG/LFP DATA
% ============================================================================

\section{S-Entropy Computation from Neural Recordings}
\label{sec:sentropy}

The S-entropy coordinates $(\Sk, \St, \Se) \in [0,1]^3$ encode the informational state of the neural system along three orthogonal axes: spatial (knowledge), temporal, and energetic entropy.
We derive their computation from multichannel EEG or local field potential (LFP) recordings.

\subsection{Setup}

Consider $N_{\mathrm{ch}}$ recording channels with sampling rate $f_s$ (typically 256~Hz for clinical EEG).
Let $x_i(t)$ denote the signal at channel $i$, $i = 1, \ldots, N_{\mathrm{ch}}$, over a window of $T$ samples.

\begin{definition}[Spatial Entropy $\Sk$]
\label{def:sk}
The spatial (knowledge) entropy measures the diversity of instantaneous amplitude distributions across channels.
For each time sample $t$, compute the normalized amplitude vector:
\begin{equation}
p_i(t) = \frac{|x_i(t)|}{\sum_{j=1}^{N_{\mathrm{ch}}} |x_j(t)| + \epsilon},
\end{equation}
where $\epsilon > 0$ prevents division by zero (Axiom~\ref{ax:nonull} ensures the denominator is generically nonzero).
The spatial Shannon entropy at time $t$ is
\begin{equation}
H_k(t) = -\sum_{i=1}^{N_{\mathrm{ch}}} p_i(t) \ln p_i(t).
\end{equation}
The S-entropy coordinate is the time-averaged, normalized value:
\begin{equation}
\boxed{\Sk = \frac{\langle H_k(t) \rangle_T}{\ln N_{\mathrm{ch}}}}
\label{eq:sk}
\end{equation}
where $\langle \cdot \rangle_T$ denotes the mean over the analysis window and $\ln N_{\mathrm{ch}}$ is the maximum entropy for $N_{\mathrm{ch}}$ channels.
\end{definition}

\begin{definition}[Temporal Entropy $\St$]
\label{def:st}
The temporal entropy measures the regularity of each channel's time series.
For channel $i$, compute the power spectral density $P_i(f)$ via FFT over the window of $T$ samples, then form the normalized spectral distribution:
\begin{equation}
q_i(f) = \frac{P_i(f)}{\sum_{f'} P_i(f')}.
\end{equation}
The spectral entropy of channel $i$ is
\begin{equation}
H_{t,i} = -\sum_f q_i(f) \ln q_i(f).
\end{equation}
The S-entropy coordinate is the channel-averaged, normalized value:
\begin{equation}
\boxed{\St = \frac{\langle H_{t,i} \rangle_{i}}{\ln N_f}}
\label{eq:st}
\end{equation}
where $N_f = T/2$ is the number of frequency bins and $\langle \cdot \rangle_i$ denotes the mean over channels.
\end{definition}

\begin{definition}[Energetic Entropy $\Se$]
\label{def:se}
The energetic entropy measures the distribution of power across frequency bands.
Partition the spectrum into $B$ standard bands (delta: 0.5--4~Hz, theta: 4--8~Hz, alpha: 8--13~Hz, beta: 13--30~Hz, gamma: 30--100~Hz; $B = 5$).
For the channel-averaged power spectrum $\bar{P}(f) = N_{\mathrm{ch}}^{-1} \sum_i P_i(f)$, compute the band power fractions:
\begin{equation}
r_b = \frac{\int_{f \in \text{band } b} \bar{P}(f)\,df}{\int_0^{f_s/2} \bar{P}(f)\,df}, \quad b = 1, \ldots, B.
\end{equation}
The energetic entropy is
\begin{equation}
\boxed{\Se = \frac{-\sum_{b=1}^B r_b \ln r_b}{\ln B}}
\label{eq:se}
\end{equation}
\end{definition}

\begin{proposition}[S-Entropy Bounds]
\label{prop:sbounds}
Under Axioms~\ref{ax:bounded}--\ref{ax:resolution}, the S-entropy coordinates satisfy $(\Sk, \St, \Se) \in [0,1]^3$ with equality at the boundaries corresponding to maximal order ($0$) and maximal disorder ($1$).
\end{proposition}

\begin{proof}
Each coordinate is the ratio of a Shannon entropy to its maximum.
Shannon entropy is non-negative (from the non-negativity of $-p \ln p$ for $p \in [0,1]$) and attains its maximum $\ln K$ for the uniform distribution over $K$ outcomes.
The ratio therefore lies in $[0,1]$.
The lower bound $0$ is attained when all probability mass concentrates on a single outcome (maximal order); the upper bound $1$ when the distribution is uniform (maximal disorder).
\end{proof}

\begin{remark}[CPU-Side Computation]
\label{rem:cpu}
The S-entropy coordinates are computed on the CPU from the raw EEG/LFP signals.
They are then uploaded to the GPU as three uniform floats: \texttt{u\_Sk}, \texttt{u\_St}, \texttt{u\_Se}.
This is the \emph{only} data transfer from CPU to GPU per observation frame.
The raw signal data never touches GPU memory.
\end{remark}


% ============================================================================
%                    SECTION 4: KURAMOTO ORDER PARAMETER FROM RECORDINGS
% ============================================================================

\section{Kuramoto Order Parameter from Multichannel Recordings}
\label{sec:kuramoto}

\subsection{Analytic Signal and Instantaneous Phase}

For each EEG channel $x_i(t)$, the instantaneous phase $\phi_i(t)$ is extracted via the Hilbert transform~\cite{gabor1946,oppenheim1999}:
\begin{equation}
z_i(t) = x_i(t) + j\,\hat{x}_i(t), \quad \phi_i(t) = \arg z_i(t),
\label{eq:hilbert}
\end{equation}
where $\hat{x}_i(t)$ is the Hilbert transform of $x_i(t)$ and $j = \sqrt{-1}$.
In practice, band-pass filtering to a frequency band of interest (e.g., alpha: 8--13~Hz) precedes the Hilbert transform to ensure a well-defined instantaneous frequency~\cite{nunez2006,buzsaki2004}.
This phase-locking approach is central to the communication-through-coherence hypothesis~\cite{fries2015,varela2001}.

\subsection{Order Parameter Computation}

\begin{definition}[Kuramoto Order Parameter]
\label{def:kuramoto}
The complex-valued Kuramoto order parameter for $N_{\mathrm{ch}}$ channels is
\begin{equation}
Z(t) = R(t)\,e^{j\Psi(t)} = \frac{1}{N_{\mathrm{ch}}} \sum_{i=1}^{N_{\mathrm{ch}}} e^{j\phi_i(t)},
\label{eq:kuramoto-Z}
\end{equation}
where $R(t) = |Z(t)| \in [0,1]$ is the synchronization magnitude and $\Psi(t)$ the mean phase~\cite{kuramoto1975,lachaux1999}.
\end{definition}

\begin{proposition}[Phase Variance Relation]
\label{prop:variance}
The circular phase variance is related to the order parameter by
\begin{equation}
\sigma^2(\phi) = -2\ln R.
\label{eq:variance-R}
\end{equation}
For $R \to 1$, $\sigma^2 \to 0$ (perfect synchrony); for $R \to 0$, $\sigma^2 \to \infty$ (complete desynchrony), capped at $2\pi^2$ by the bounded phase space.
\end{proposition}

\begin{proof}
The circular variance is defined as $\mathrm{Var}_{\mathrm{circ}} = 1 - R$~\cite{acebron2005}.
For the von Mises distribution with concentration parameter $\kappa$, $R = I_1(\kappa)/I_0(\kappa)$ and $\sigma^2 = 1/\kappa$ for large $\kappa$.
The relation $\sigma^2 = -2\ln R$ follows from the small-variance expansion $R \approx \exp(-\sigma^2/2)$, which is exact for Gaussian-distributed phases~\cite{strogatz2000}.
On the bounded configuration manifold $\mathcal{M}$, we impose $\sigma^2 \leq 2\pi^2$ by Axiom~\ref{ax:bounded}.
\end{proof}

\begin{remark}[Window-Averaged Values]
For GPU upload, both $R$ and $\sigma^2$ are averaged over the analysis window:
\begin{equation}
\bar{R} = \langle R(t) \rangle_T, \quad \bar{\sigma}^2 = -2\ln\bar{R}.
\end{equation}
These are uploaded as uniforms \texttt{u\_R} and \texttt{u\_sigma2} alongside the S-entropy coordinates.
\end{remark}


% ============================================================================
%                    SECTION 5: NEURAL DOMAIN SHADER
% ============================================================================

\section{The Neural Domain Shader}
\label{sec:shader}

\subsection{Architecture}

The universal shader pipeline~\cite{sachikonye2026gpu} comprises five stages: vertex shader (fullscreen triangle), partition observation shader, interference shader, quality metric shader, and display shader.
Only the partition observation shader's \texttt{computePartitionValue()} function changes between domains.
We now specify the neural domain implementation.

\begin{theorem}[Neural Partition Observation]
\label{thm:neural-observation}
Let $F_{\mathrm{neural}} : [0,1]^2 \times \R^p \to [0,1]^4$ be the neural domain fragment shader, with texel coordinates $(u,v) \in [0,1]^2$ and parameter vector including $(\Sk, \St, \Se, R, \sigma^2)$ plus 16 domain parameters.
The output RGBA encodes:
\begin{align}
\mathrm{R} &= A(u,v;\Svec,R) \quad &&\text{(amplitude)} \label{eq:rgba-r} \\
\mathrm{G} &= \varphi(u,v;\Svec,\sigma^2) \quad &&\text{(phase)} \label{eq:rgba-g} \\
\mathrm{B} &= n(u,v;\Svec) \quad &&\text{(partition depth)} \label{eq:rgba-b} \\
\mathrm{A} &= \kappa(u,v;\Svec,R) \quad &&\text{(confidence)} \label{eq:rgba-a}
\end{align}
where $A$ is the partition amplitude, $\varphi$ the local oscillatory phase, $n$ the principal quantum number (shell depth), and $\kappa$ the observation confidence.
By the Rendering--Measurement Identity~\cite{sachikonye2026gpu}, evaluating $F_{\mathrm{neural}}$ at every texel is identical to performing a complete categorical measurement of the neural partition state.
\end{theorem}

\begin{proof}
The Rendering--Measurement Identity requires four conditions~\cite{sachikonye2026gpu}: (i)~input identity (texel coordinates map to partition cell labels), (ii)~per-texel operation identity (the shader evaluates a categorical observable), (iii)~parallelism identity (GPU independence maps to measurement independence via the commutation theorem), (iv)~output identity (the RGBA vector encodes a complete categorical record).

For the neural domain: (i)~the texel coordinate $(u,v)$ is mapped to partition coordinates $(n,\ell) = (\lceil u \cdot N \rceil, \lfloor v \cdot n \rfloor)$ where $N$ is the maximum shell; (ii)~the function \texttt{computePartitionValue()} evaluates the amplitude and phase of the neural oscillatory mode at shell $n$, sector $\ell$, with coupling determined by $R$ and variance $\sigma^2$; (iii)~each texel is independent since the observation of partition cell $(n,\ell)$ commutes with observation of cell $(n',\ell')$; (iv)~the four RGBA channels provide a complete record of amplitude, phase, depth, and confidence.
All four conditions hold; the identity follows.
\end{proof}

\subsection{GLSL Implementation}

The following GLSL~ES~3.00~\cite{khronos2023glsl,khronos2023webgl2} code implements the neural domain partition observation shader.

\begin{lstlisting}[style=glsl,caption={Neural domain \texttt{computePartitionValue()} --- the fragment shader function that evaluates the neural partition state at each texel. Uniforms are uploaded from the CPU-computed S-entropy and Kuramoto values.},label=lst:partition]
#version 300 es
precision highp float;

// S-entropy coordinates from CPU
uniform float u_Sk;   // spatial entropy
uniform float u_St;   // temporal entropy
uniform float u_Se;   // energetic entropy
// Kuramoto parameters from CPU
uniform float u_R;    // order parameter
uniform float u_sigma2; // phase variance
// Domain parameters
uniform float u_domainParams[16];
// u_domainParams[0] = max shell N
// u_domainParams[1] = coupling K
// u_domainParams[2] = critical K_c
// u_domainParams[3] = temperature T
// u_domainParams[4] = time t
// u_domainParams[5..15] reserved

uniform vec2 u_resolution;

out vec4 fragColor;

const float PI = 3.14159265359;
const float TWO_PI = 6.28318530718;

// Shell capacity: C(n) = 2n^2
float shellCapacity(float n) {
  return 2.0 * n * n;
}

// Cumulative capacity up to shell N
float cumulativeCapacity(float N) {
  return N * (N + 1.0) * (2.0 * N + 1.0)
         / 3.0;
}

// Structural factor S = R * exp(-s2/2pi^2)
float structuralFactor(float R, float s2) {
  return R * exp(-s2 / (2.0 * PI * PI));
}

// Regime classification from R
int classifyRegime(float R) {
  if (R < 0.3) return 0;      // turbulent
  if (R < 0.5) return 1;      // aperture
  if (R < 0.8) return 2;      // cascade
  if (R < 0.95) return 3;     // coherent
  return 4;                    // phase-locked
}

// Partition amplitude at shell n, sector l
// Derived from Triple Equivalence:
//   amplitude ~ sqrt(C(n)/C_tot) * entropy
float partitionAmplitude(
  float n, float l, float Sk, float St,
  float Se, float R
) {
  float maxN = u_domainParams[0];
  float Cn = shellCapacity(n);
  float Ctot = cumulativeCapacity(maxN);
  // Radial weight from shell occupation
  float radialW = sqrt(Cn / Ctot);
  // Angular weight from sector within shell
  float angularW = (2.0 * l + 1.0)
                   / (n * n);
  // S-entropy modulation: partition density
  // concentrates where entropy is high
  float entropyMod = (Sk + St + Se) / 3.0;
  // Synchronization amplification
  float syncAmp = mix(0.1, 1.0, R);
  return radialW * angularW
         * entropyMod * syncAmp;
}

// Partition phase from oscillatory mode
float partitionPhase(
  float n, float l, float sigma2, float t
) {
  // Each (n,l) mode has natural frequency
  // omega_{n,l} = 2*pi*(n - l/(n+1))
  float omega = TWO_PI
    * (n - l / (n + 1.0));
  // Phase noise from variance
  float noise = sqrt(sigma2)
    * sin(omega * 0.1 + n * l);
  return mod(omega * t + noise, TWO_PI)
         / TWO_PI; // normalize to [0,1]
}

// Main partition value computation
vec4 computePartitionValue(vec2 uv) {
  float maxN = u_domainParams[0];
  float t = u_domainParams[4];

  // Map texel to partition coordinates
  // u -> principal shell n in [1, maxN]
  float n = floor(uv.x * maxN) + 1.0;
  n = clamp(n, 1.0, maxN);
  // v -> angular sector l in [0, n-1]
  float l = floor(uv.y * n);
  l = clamp(l, 0.0, n - 1.0);

  // Compute RGBA observation
  float amp = partitionAmplitude(
    n, l, u_Sk, u_St, u_Se, u_R);
  float phase = partitionPhase(
    n, l, u_sigma2, t);
  float depth = n / maxN; // normalized
  // Confidence: higher at lower variance,
  //   scaled by structural factor
  float S = structuralFactor(u_R, u_sigma2);
  float conf = S * (1.0 - depth * 0.3);

  return vec4(amp, phase, depth, conf);
}

void main() {
  vec2 uv = gl_FragCoord.xy / u_resolution;
  fragColor = computePartitionValue(uv);
}
\end{lstlisting}

\begin{remark}[Per-Texel Independence]
Each texel's computation depends only on the uniform parameters and its own $(u,v)$ coordinates.
No texel reads from any other texel.
This guarantees the parallelism condition of the Rendering--Measurement Identity and ensures that the GPU's SIMD execution model produces identical results regardless of execution order.
\end{remark}


% ============================================================================
%                    SECTION 6: STRUCTURAL FACTOR COMPUTATION
% ============================================================================

\section{Structural Factor on GPU}
\label{sec:structural}

\begin{definition}[Structural Factor]
\label{def:structural}
The structural factor combines synchronization strength with phase coherence~\cite{sachikonye2026regime,sachikonye2026lagrangian}:
\begin{equation}
\boxed{S = R \exp\!\left(-\frac{\sigma^2}{2\pi^2}\right)}
\label{eq:structural}
\end{equation}
with $S \in [0,1]$, where $S = 1$ corresponds to perfect phase-locking ($R = 1$, $\sigma^2 = 0$) and $S \to 0$ corresponds to either weak synchronization or large phase variance.
\end{definition}

\begin{theorem}[Per-Texel Structural Factor]
\label{thm:structural-texel}
The structural factor $S$ can be evaluated at each texel as a function of the uniform parameters $(R, \sigma^2)$ without inter-texel communication.
The per-texel evaluation cost is exactly three floating-point operations: one negation, one division, one \texttt{exp}, and one multiplication.
\end{theorem}

\begin{proof}
The structural factor $S = R \exp(-\sigma^2 / 2\pi^2)$ depends only on $R$ and $\sigma^2$, both of which are uniform parameters broadcast to all texels.
No texel-local data enters the computation of $S$ itself.
The GPU evaluates $\texttt{exp}$ in hardware (single-cycle on modern GPUs), so the cost is:
(i)~compute $-\sigma^2/(2\pi^2)$: one negation, one division (or multiplication by the precomputed constant $1/(2\pi^2)$);
(ii)~evaluate $\exp$: one hardware instruction;
(iii)~multiply by $R$: one multiplication.
Total: $\BigO(1)$ per texel with negligible constant.
\end{proof}

\begin{proposition}[Structural Factor as Regime Discriminant]
\label{prop:s-discriminant}
The structural factor $S$ provides a single scalar that discriminates between regimes more sharply than $R$ alone.
Specifically, using~\eqref{eq:variance-R}:
\begin{equation}
S = R \exp\!\left(\frac{\ln R}{\pi^2}\right) = R^{1 + 1/\pi^2}.
\label{eq:s-from-R}
\end{equation}
The exponent $1 + 1/\pi^2 \approx 1.101$ amplifies differences near regime boundaries.
\end{proposition}

\begin{proof}
Substituting $\sigma^2 = -2\ln R$ from Proposition~\ref{prop:variance} into~\eqref{eq:structural}:
\begin{equation}
S = R \exp\!\left(\frac{-(-2\ln R)}{2\pi^2}\right) = R \exp\!\left(\frac{\ln R}{\pi^2}\right) = R \cdot R^{1/\pi^2} = R^{1+1/\pi^2}.
\end{equation}
Since $1 + 1/\pi^2 > 1$, the map $R \mapsto S$ is a power law with exponent greater than unity, which steepens the transition at intermediate $R$ values (precisely the aperture and cascade regimes where discrimination is most needed).
\end{proof}


% ============================================================================
%                    SECTION 7: EULER--LAGRANGE INTEGRATION ON GPU
% ============================================================================

\section{Euler--Lagrange Gradient Flow on GPU}
\label{sec:euler-lagrange}

\subsection{The Neural Partition Lagrangian}

The Onsager--Machlup Lagrangian on the generalized coordinates $\qvec = (R, \sigma^2, \Sk, \St, \Se)$ is~\cite{sachikonye2026lagrangian,onsager1953,machlup1953}:
\begin{equation}
\Lag(\qvec, \dot{\qvec}) = \frac{1}{2} g_{ij}\,\dot{q}^i\,\dot{q}^j - \Pot(\qvec),
\label{eq:lagrangian}
\end{equation}
where $g_{ij}$ is the metric tensor on $\mathcal{M}$ and $\Pot(\qvec)$ is the neural partition potential.

\begin{definition}[Neural Partition Potential]
\label{def:potential}
The potential decomposes into four terms:
\begin{equation}
\boxed{\Pot(\qvec) = V_{\mathrm{sync}} + V_{\mathrm{var}} + V_{\mathrm{SF}} + V_{\mathrm{ent}}}
\label{eq:potential}
\end{equation}
with:
\begin{align}
V_{\mathrm{sync}} &= -\frac{a}{2}R^2 + \frac{b}{4}R^4, \quad a = K - K_c, \label{eq:vsync} \\
V_{\mathrm{var}} &= \frac{\lambda}{2}\!\left(\sigma^2 - \frac{\kb T}{K}\right)^{\!2}, \label{eq:vvar} \\
V_{\mathrm{SF}} &= -\mu\,R\exp\!\left(-\frac{\sigma^2}{2\pi^2}\right), \label{eq:vsf} \\
V_{\mathrm{ent}} &= \nu \sum_{\alpha \in \{k,t,e\}} \left[S_\alpha \ln S_\alpha + (1-S_\alpha)\ln(1-S_\alpha)\right]. \label{eq:vent}
\end{align}
Here $K$ is the coupling strength, $K_c = 2\sigma_\omega/\pi$ the critical coupling~\cite{kuramoto1975}, $\lambda, \mu, \nu > 0$ are stiffness constants, and $T$ the temperature.
\end{definition}

\subsection{Euler--Lagrange Equations}

In the overdamped (gradient flow) limit appropriate for neural dynamics, the Euler--Lagrange equations reduce to~\cite{sachikonye2026lagrangian}:
\begin{align}
\dot{R} &= -\frac{\partial \Pot}{\partial R} = aR - bR^3 + \mu\exp\!\left(-\frac{\sigma^2}{2\pi^2}\right), \label{eq:rdot} \\
\dot{\sigma}^2 &= -\frac{\partial \Pot}{\partial \sigma^2} = -\lambda\!\left(\sigma^2 - \frac{\kb T}{K}\right) - \frac{\mu R}{2\pi^2}\exp\!\left(-\frac{\sigma^2}{2\pi^2}\right), \label{eq:s2dot} \\
\dot{S}_\alpha &= -\frac{\partial \Pot}{\partial S_\alpha} = -\nu \ln\!\frac{S_\alpha}{1 - S_\alpha}, \quad \alpha \in \{k,t,e\}. \label{eq:sdot}
\end{align}

\begin{theorem}[Fixed Points and Regime Correspondence]
\label{thm:fixedpoints}
The fixed points $\dot{\qvec} = 0$ of the gradient flow~\eqref{eq:rdot}--\eqref{eq:sdot} correspond to the five operational regimes.
Specifically:
\begin{enumerate}[label=(\roman*),leftmargin=*]
\item For $K < K_c$ ($a < 0$): $R^* \approx 0$, turbulent regime.
\item For $K$ slightly above $K_c$: $R^* = \sqrt{a/b}$ small, aperture-dominated regime.
\item For intermediate $K$: $R^*$ in the cascade range, with $S_\alpha^* = 1/2$ (maximum entropy) at the entropic fixed point.
\item For large $K$: $R^*$ approaches unity, coherent regime.
\item For $K \gg K_c$: $R^* \to 1$, $\sigma^2 \to 0$, phase-locked regime.
\end{enumerate}
\end{theorem}

\begin{proof}
Setting $\dot{R} = 0$ in~\eqref{eq:rdot} with the structural factor term small relative to the Landau terms gives $R^*(aR - bR^3) \approx 0$, yielding $R^* = 0$ or $R^* = \sqrt{a/b}$.
The pitchfork bifurcation at $a = 0$ (i.e., $K = K_c$) separates the turbulent ($R^* = 0$) from the ordered ($R^* > 0$) regimes~\cite{strogatz2000}.
The S-entropy equation~\eqref{eq:sdot} has $\dot{S}_\alpha = 0$ when $S_\alpha / (1 - S_\alpha) = 1$, i.e., $S_\alpha^* = 1/2$.
The regime classification follows from mapping the fixed-point value of $R^*$ to Definition~\ref{def:regimes}.
\end{proof}

\subsection{GPU Integration Shader}

The following shader integrates the Euler--Lagrange equations for one time step using the explicit Euler method.
The state vector is stored in a texture and updated each frame.

\begin{lstlisting}[style=glsl,caption={Euler--Lagrange gradient flow integration shader. Reads the current state from a texture, computes the time derivatives from the neural partition potential, and writes the updated state. One integration step per draw call.},label=lst:euler]
#version 300 es
precision highp float;

uniform sampler2D u_stateTex; // current q
uniform float u_dt;           // time step
uniform float u_K;            // coupling
uniform float u_Kc;           // critical K
uniform float u_kBT;          // k_B * T
uniform float u_lambda;       // var stiffness
uniform float u_mu;           // SF coupling
uniform float u_nu;           // ent stiffness
uniform float u_b;            // quartic coeff
uniform vec2 u_resolution;

out vec4 fragColor;

const float PI = 3.14159265359;
const float TWO_PI_SQ = 2.0 * PI * PI;

void main() {
  vec2 uv = gl_FragCoord.xy / u_resolution;

  // Read current state:
  // texel (0,0): (R, sigma2, Sk, St)
  // texel (1,0): (Se, 0, 0, 0)
  vec4 s0 = texture(u_stateTex,
    vec2(0.5, 0.5) / u_resolution);
  vec4 s1 = texture(u_stateTex,
    vec2(1.5, 0.5) / u_resolution);

  float R = s0.x;
  float sigma2 = s0.y;
  float Sk = s0.z;
  float St = s0.w;
  float Se = s1.x;

  // Landau parameter
  float a = u_K - u_Kc;

  // Structural factor exponential
  float expTerm = exp(-sigma2 / TWO_PI_SQ);

  // dR/dt = aR - bR^3 + mu*exp(...)
  float dR = a * R - u_b * R * R * R
           + u_mu * expTerm;

  // dsigma2/dt = -lambda*(sigma2 - kBT/K)
  //   - mu*R/(2*pi^2)*exp(...)
  float varFloor = u_kBT / max(u_K, 1e-6);
  float dSigma2 = -u_lambda * (sigma2 - varFloor)
    - u_mu * R / TWO_PI_SQ * expTerm;

  // dSk/dt = -nu * ln(Sk/(1-Sk))
  float dSk = -u_nu
    * log(max(Sk, 1e-6)
        / max(1.0 - Sk, 1e-6));
  float dSt = -u_nu
    * log(max(St, 1e-6)
        / max(1.0 - St, 1e-6));
  float dSe = -u_nu
    * log(max(Se, 1e-6)
        / max(1.0 - Se, 1e-6));

  // Euler integration
  R = clamp(R + dR * u_dt, 0.0, 1.0);
  sigma2 = clamp(sigma2 + dSigma2 * u_dt,
                 0.0, TWO_PI_SQ);
  Sk = clamp(Sk + dSk * u_dt, 0.001, 0.999);
  St = clamp(St + dSt * u_dt, 0.001, 0.999);
  Se = clamp(Se + dSe * u_dt, 0.001, 0.999);

  // Write updated state
  if (gl_FragCoord.x < 1.0) {
    fragColor = vec4(R, sigma2, Sk, St);
  } else {
    fragColor = vec4(Se, 0.0, 0.0, 1.0);
  }
}
\end{lstlisting}

\begin{remark}[Integration Stability]
The explicit Euler method is stable for the gradient flow provided $\Delta t < 2/\lambda_{\max}$, where $\lambda_{\max}$ is the largest eigenvalue of the Hessian $\partial^2 \Pot / \partial q^i \partial q^j$.
For typical neural parameters ($\lambda \sim 1$, $\mu \sim 0.1$, $\nu \sim 0.5$), $\Delta t = 0.01$ suffices.
At 60~fps, each frame advances the trajectory by $\Delta t = 1/60 \approx 0.017$, within the stability bound.
Higher-order methods (RK4) can be implemented as multi-pass shaders if needed.
\end{remark}


% ============================================================================
%                    SECTION 8: INTERFERENCE-BASED REGIME MATCHING
% ============================================================================

\section{Interference-Based Neural State Matching}
\label{sec:interference}

\subsection{The Universal Interference Shader}

The universal interference shader~\cite{sachikonye2026gpu} superposes two observation textures as wave amplitudes:

\begin{definition}[Interference Operation]
\label{def:interference}
Given two observation textures $\Tex_q$ (query) and $\Tex_r$ (reference), with per-texel amplitude $A$ and phase $\varphi$ stored in channels R and G, the interference intensity at each texel is:
\begin{equation}
\boxed{I(u,v) = \left|A_q\,e^{i\varphi_q} + A_r\,e^{i\varphi_r}\right|^2 = A_q^2 + A_r^2 + 2A_q A_r \cos(\varphi_q - \varphi_r)}
\label{eq:interference}
\end{equation}
\end{definition}

\begin{definition}[Interference Visibility]
\label{def:visibility}
The visibility $\mathcal{V}$ is computed from the max and min of the interference pattern:
\begin{equation}
\mathcal{V} = \frac{I_{\max} - I_{\min}}{I_{\max} + I_{\min}}.
\label{eq:visibility}
\end{equation}
For identical states $\Tex_q = \Tex_r$, constructive interference gives $\mathcal{V} = 1$.
For orthogonal states, $\mathcal{V} = 0$.
\end{definition}

\begin{theorem}[Regime Identification by Interference]
\label{thm:regime-interference}
Let $\Tex_q$ be the observation texture of an unknown neural state and $\{\Tex_r^{(k)}\}_{k=0}^{4}$ the five regime signature textures (computed once from canonical parameter values for each regime).
The regime of the unknown state is:
\begin{equation}
k^* = \arg\max_k \mathcal{V}(\Tex_q, \Tex_r^{(k)}).
\label{eq:regime-id}
\end{equation}
If $\mathcal{V}(\Tex_q, \Tex_r^{(k^*)}) \geq 0.95$, the identification is definitive.
\end{theorem}

\begin{proof}
The observation texture encodes the complete categorical measurement of the neural partition state (Theorem~\ref{thm:neural-observation}).
By the Texture-State Isomorphism~\cite{sachikonye2026gpu}, the $L^2$ distance between textures is proportional to the categorical distance between states:
$\dcat(\sigma_1, \sigma_2) = \lambda \|\Tex_1 - \Tex_2\|_{L^2}$.
Interference visibility is maximized when the two textures are identical (constructive interference at all texels), corresponding to zero categorical distance.
Since the five regimes occupy disjoint regions of $\mathcal{M}$ (Definition~\ref{def:regimes}), the canonical regime textures are mutually distinguishable, and the argmax over $k$ uniquely identifies the regime.
The threshold $\mathcal{V} \geq 0.95$ accounts for noise and finite-resolution effects.
\end{proof}

\subsection{Regime Signature Textures}

The canonical regime signatures are precomputed once from the following parameter values:

\begin{center}
\begin{tabular}{lccccc}
\toprule
Regime & $R$ & $\sigma^2$ & $\Sk$ & $\St$ & $\Se$ \\
\midrule
Turbulent & 0.15 & 3.79 & 0.90 & 0.85 & 0.80 \\
Aperture & 0.40 & 1.83 & 0.70 & 0.65 & 0.60 \\
Cascade & 0.65 & 0.86 & 0.50 & 0.50 & 0.50 \\
Coherent & 0.88 & 0.26 & 0.30 & 0.35 & 0.40 \\
Phase-locked & 0.97 & 0.06 & 0.15 & 0.20 & 0.25 \\
\bottomrule
\end{tabular}
\end{center}

\noindent The $\sigma^2$ values are computed from $R$ via~\eqref{eq:variance-R}.
These five textures require $5 \times 512 \times 512 \times 4 \times 4 = 20$~MB at float32 precision, or $5$~MB at RGBA8.

% --- Figure 1 placeholder ---
\begin{figure}[t]
\centering
\fbox{\parbox{0.9\columnwidth}{\centering\vspace{2cm}
\textbf{Figure 1:} Interference visibility matrix $\mathcal{V}(k,k')$ between the five canonical regime textures. Diagonal entries show $\mathcal{V} = 1.0$ (self-matching); off-diagonal entries decrease with regime distance, confirming mutual distinguishability. Adjacent regimes show $\mathcal{V} \approx 0.3$--$0.5$; distant regimes show $\mathcal{V} < 0.1$.
\vspace{2cm}}}
\caption{Regime interference visibility matrix.}
\label{fig:visibility-matrix}
\end{figure}


% ============================================================================
%                    SECTION 9: GPU PHYSICAL OBSERVABLES
% ============================================================================

\section{GPU Physical Observables for Neural State Classification}
\label{sec:observables}

The universal quality metric shader extracts four physical observables from the observation texture.
These provide objective, deterministic quality metrics for neural state classification without human labels.

\begin{definition}[Partition Sharpness]
\label{def:sharpness}
The partition sharpness at texel $(u,v)$ is the gradient magnitude of the amplitude channel:
\begin{equation}
\mathcal{P}(u,v) = \|\nabla A(u,v)\| = \sqrt{\left(\frac{\partial A}{\partial u}\right)^2 + \left(\frac{\partial A}{\partial v}\right)^2}.
\label{eq:sharpness}
\end{equation}
High sharpness indicates well-resolved partition boundaries; low sharpness indicates a blurred, degenerate state.
\end{definition}

\begin{definition}[Phase Coherence]
\label{def:phase-coherence}
The local phase coherence at texel $(u,v)$ is computed from the four nearest neighbors:
\begin{equation}
\mathcal{C}(u,v) = \left|\frac{1}{4}\sum_{(u',v') \in \mathcal{N}} e^{2\pi i\,\varphi(u',v')}\right|,
\label{eq:coherence}
\end{equation}
where $\mathcal{N}$ is the 4-connected neighborhood.
$\mathcal{C} = 1$ indicates perfect local phase agreement; $\mathcal{C} = 0$ indicates complete phase disorder.
\end{definition}

\begin{definition}[Multi-Resolution Consistency]
\label{def:multiresolution}
The multi-resolution consistency $\mathcal{R}$ measures agreement between the observation texture at resolution $W \times W$ and a $2\times$-downsampled version at $(W/2) \times (W/2)$:
\begin{equation}
\mathcal{R} = 1 - \frac{\|\Tex_W - \mathrm{Up}(\Tex_{W/2})\|_{L^2}}{\|\Tex_W\|_{L^2} + \epsilon},
\label{eq:multiresolution}
\end{equation}
where $\mathrm{Up}$ denotes bilinear upsampling.
Legitimate partition structure is scale-invariant ($\mathcal{R} \to 1$); noise is not ($\mathcal{R} \to 0$).
\end{definition}

\begin{theorem}[Observable-Regime Mapping]
\label{thm:observable-regime}
The four GPU observables $(\mathcal{P}, \mathcal{C}, \mathcal{V}, \mathcal{R})$ provide a bijective mapping to the five operational regimes:
\begin{center}
\begin{tabular}{lcccc}
\toprule
Regime & $\mathcal{P}$ & $\mathcal{C}$ & $\mathcal{V}$ & $\mathcal{R}$ \\
\midrule
Turbulent & low & low & low & low \\
Aperture & med & low--med & med & med \\
Cascade & high & med & med--high & high \\
Coherent & med--high & high & high & high \\
Phase-locked & low & $\to 1$ & $\to 1$ & $\to 1$ \\
\bottomrule
\end{tabular}
\end{center}
\end{theorem}

\begin{proof}
In the turbulent regime ($R < 0.3$), the partition amplitudes are weakly modulated (low $\mathcal{P}$), phases are random (low $\mathcal{C}$), and downsampling destroys what little structure exists (low $\mathcal{R}$).
In the aperture-dominated regime, selective gating creates moderate amplitude gradients at aperture boundaries (med $\mathcal{P}$) with partial phase ordering.
In the hierarchical cascade, the multi-scale structure produces sharp partition boundaries (high $\mathcal{P}$) that are scale-invariant (high $\mathcal{R}$).
In the coherent regime, global synchronization produces high phase coherence ($\mathcal{C} \to 1$) with smooth amplitude variation (med--high $\mathcal{P}$).
In the phase-locked regime, all phases align ($\mathcal{C} = 1$), the observation texture becomes nearly uniform (low $\mathcal{P}$), and scale invariance is perfect ($\mathcal{R} = 1$).
The ordering is non-degenerate: no two regimes share the same profile across all four observables.
\end{proof}


% ============================================================================
%                    SECTION 10: REGIME CLASSIFICATION SHADER
% ============================================================================

\section{Regime Classification Logic}
\label{sec:regime-shader}

The following shader implements regime classification from the GPU observables and the structural factor.

\begin{lstlisting}[style=glsl,caption={Regime classification shader. Reads the quality metric texture and uniform parameters, outputs the classified regime as a color-coded value and the structural factor. Combines both threshold-based (from $R$) and observable-based classification.},label=lst:regime]
#version 300 es
precision highp float;

uniform sampler2D u_qualityTex;
uniform float u_R;
uniform float u_sigma2;
uniform vec2 u_resolution;

out vec4 fragColor;

const float PI = 3.14159265359;

// Structural factor
float structuralFactor(float R, float s2) {
  return R * exp(-s2 / (2.0 * PI * PI));
}

// Regime boundaries
const float TURB_BOUND  = 0.3;
const float APER_BOUND  = 0.5;
const float CASC_BOUND  = 0.8;
const float COHER_BOUND = 0.95;

// Regime colors for display
const vec3 COL_TURB  = vec3(0.2, 0.2, 0.8);
const vec3 COL_APER  = vec3(0.2, 0.7, 0.3);
const vec3 COL_CASC  = vec3(0.9, 0.8, 0.1);
const vec3 COL_COHER = vec3(0.9, 0.4, 0.1);
const vec3 COL_LOCK  = vec3(0.9, 0.1, 0.1);

void main() {
  vec2 uv = gl_FragCoord.xy / u_resolution;

  // Read quality metrics from texture
  // R=sharpness, G=coherence,
  // B=visibility, A=consistency
  vec4 q = texture(u_qualityTex, uv);
  float sharpness   = q.r;
  float coherence   = q.g;
  float visibility  = q.b;
  float consistency = q.a;

  // Compute structural factor
  float S = structuralFactor(u_R, u_sigma2);

  // Primary classification from R
  int regime = 0;
  vec3 regimeColor = COL_TURB;
  if (u_R >= COHER_BOUND) {
    regime = 4;
    regimeColor = COL_LOCK;
  } else if (u_R >= CASC_BOUND) {
    regime = 3;
    regimeColor = COL_COHER;
  } else if (u_R >= APER_BOUND) {
    regime = 2;
    regimeColor = COL_CASC;
  } else if (u_R >= TURB_BOUND) {
    regime = 1;
    regimeColor = COL_APER;
  }

  // Confidence from observables agreement
  float conf = 0.0;
  if (regime == 0) { // turbulent
    conf = (1.0 - coherence) * (1.0 - S);
  } else if (regime == 4) { // phase-locked
    conf = coherence * S;
  } else { // intermediate
    conf = consistency * S;
  }
  conf = clamp(conf, 0.0, 1.0);

  // Output: RGB=regime color, A=confidence
  fragColor = vec4(regimeColor * conf, S);
}
\end{lstlisting}


% ============================================================================
%                    SECTION 11: OPERATOR TRAJECTORIES
% ============================================================================

\section{Operator Trajectories Through Regime Space}
\label{sec:trajectories}

\begin{definition}[Operator Trajectory]
\label{def:trajectory}
An operator trajectory is a continuous path $\gamma : [0, T] \to \mathcal{M}$ through the configuration manifold, with $\gamma(t) = (R(t), \sigma^2(t), \Sk(t), \St(t), \Se(t))$.
The trajectory satisfies the Euler--Lagrange gradient flow~\eqref{eq:rdot}--\eqref{eq:sdot}.
\end{definition}

\begin{definition}[Trajectory-Terminus-Memory Triple]
\label{def:ttm}
A complete system descriptor is the triple~\cite{sachikonye2026operator}
\begin{equation}
\mathcal{M} = (\gamma, \Gamma_f, M),
\label{eq:ttm}
\end{equation}
where $\gamma$ is the trajectory, $\Gamma_f = \gamma(T)$ the terminus (final state), and $M$ the memory (the observation textures accumulated along the path).
\end{definition}

\begin{theorem}[Trajectory Non-Identity]
\label{thm:nonidentity}
Let $\gamma_1$ and $\gamma_2$ be two trajectories with identical terminus $\Gamma_f = \gamma_1(T_1) = \gamma_2(T_2)$.
If $\gamma_1 \neq \gamma_2$ as paths (differing in at least one point), then:
\begin{equation}
(\gamma_1, \Gamma_f, M_1) \neq (\gamma_2, \Gamma_f, M_2).
\label{eq:nonidentity}
\end{equation}
Different paths to the same terminus yield different system states.
\end{theorem}

\begin{proof}
The memory $M$ is the sequence of observation textures $\{\Tex_{\gamma(t_k)}\}_{k=0}^{K}$ sampled along the trajectory.
Since $\gamma_1 \neq \gamma_2$, there exists $t^*$ such that $\gamma_1(t^*) \neq \gamma_2(t^*)$.
By Theorem~\ref{thm:neural-observation}, the observation textures $\Tex_{\gamma_1(t^*)} \neq \Tex_{\gamma_2(t^*)}$ (different partition states produce different textures).
Therefore $M_1 \neq M_2$, and the triples are distinct.
\end{proof}

\begin{corollary}[Drug Trajectory Specificity]
Two drugs that achieve the same final regime (e.g., both reaching the coherent regime) through different pharmacological mechanisms trace different trajectories and thus produce different therapeutic outcomes.
This is observable on GPU as different sequences of observation textures along the path.
\end{corollary}

% --- Figure 2 placeholder ---
\begin{figure}[t]
\centering
\fbox{\parbox{0.9\columnwidth}{\centering\vspace{2.5cm}
\textbf{Figure 2:} Two operator trajectories from the turbulent regime to the coherent regime, corresponding to different pharmacological interventions. Trajectory $\gamma_1$ (SSRI, blue) passes through the aperture-dominated regime; $\gamma_2$ (SNRI, red) passes through the hierarchical cascade. Both reach $R \approx 0.88$ but the accumulated memory textures differ, consistent with Theorem~\ref{thm:nonidentity}. The structural factor $S(t)$ is plotted alongside.
\vspace{2.5cm}}}
\caption{Drug-induced regime traversal: trajectory non-identity.}
\label{fig:trajectories}
\end{figure}


% ============================================================================
%                    SECTION 12: POINCARE COMPUTING ON GPU
% ============================================================================

\section{Poincar\'{e} Computing on GPU}
\label{sec:poincare}

\subsection{Backward Trajectory Integration}

Poincar\'{e} Computing~\cite{sachikonye2026operator,poincare1890} inverts the standard computational paradigm: instead of specifying initial conditions and integrating forward, one specifies the terminus (desired final state) and derives the trajectory backward.

\begin{definition}[Poincar\'{e} Problem]
\label{def:poincare}
Given a target terminus $\Gamma_f \in \mathcal{M}$ and the Euler--Lagrange equations~\eqref{eq:rdot}--\eqref{eq:sdot}, find the trajectory $\gamma : [-T, 0] \to \mathcal{M}$ such that $\gamma(0) = \Gamma_f$ and $\gamma(-T) = \Gamma_0$ lies in a specified initial regime.
\end{definition}

\begin{theorem}[Backward Integration on GPU]
\label{thm:backward}
The backward trajectory from terminus $\Gamma_f$ is computed by integrating the time-reversed gradient flow:
\begin{equation}
\frac{d\qvec}{d\tau} = +\frac{\partial \Pot}{\partial \qvec}\bigg|_{\qvec = \gamma(\tau)}, \quad \tau = -t,
\label{eq:backward}
\end{equation}
where $\tau$ is the reversed time.
The GPU integration shader (Listing~\ref{lst:euler}) is reused with $\Delta t \to -\Delta t$.
\end{theorem}

\begin{proof}
The gradient flow $\dot{\qvec} = -\partial\Pot/\partial\qvec$ under the substitution $\tau = -t$ becomes $d\qvec/d\tau = +\partial\Pot/\partial\qvec$, which is the \emph{gradient ascent} of $\Pot$.
Starting from the terminus $\Gamma_f$ (a local minimum of $\Pot$ within the target regime), gradient ascent climbs out of the potential basin, tracing the reverse of the most probable forward trajectory (by the Onsager--Machlup principle~\cite{onsager1953}).
On GPU, this requires only negating the time step sign in Listing~\ref{lst:euler}.
\end{proof}

\begin{corollary}[Treatment Planning]
Given a desired therapeutic terminus (e.g., the coherent regime with $R^* = 0.88$, $\sigma^2 = 0.26$), Poincar\'{e} backward integration computes the trajectory to that terminus.
The trajectory's passage through intermediate regimes constrains the pharmacological mechanism required.
\end{corollary}

% --- Figure 3 placeholder ---
\begin{figure}[t]
\centering
\fbox{\parbox{0.9\columnwidth}{\centering\vspace{2.5cm}
\textbf{Figure 3:} Poincar\'{e} backward computation from a coherent-regime terminus. The reversed trajectory (dashed) climbs the potential landscape from the coherent-regime basin through the cascade and aperture regimes to the turbulent initial state. The forward trajectory (solid) follows the same path in the forward direction. Observation textures at four points along the trajectory are shown as insets.
\vspace{2.5cm}}}
\caption{Poincar\'{e} Computing: backward trajectory from terminus.}
\label{fig:poincare}
\end{figure}


% ============================================================================
%                    SECTION 13: APPLICATIONS
% ============================================================================

\section{Applications}
\label{sec:applications}

\subsection{Real-Time EEG Regime Classification}

For an EEG recording with $N_{\mathrm{ch}} = 19$ channels (10--20 system) at $f_s = 256$~Hz:

\begin{enumerate}[leftmargin=*,itemsep=2pt]
\item \textbf{CPU:} Compute $(\Sk, \St, \Se, R, \sigma^2)$ from a 2-second window (512 samples).
The Hilbert transform for $R$ requires $\BigO(N_{\mathrm{ch}} \cdot T \log T) \approx 19 \times 512 \times 9 \approx 87{,}552$ operations per window.
At 1~GHz, this takes $<0.1$~ms.

\item \textbf{GPU:} Upload five uniforms.
Execute the partition observation shader on a $512 \times 512$ texture (one draw call, $\sim\!0.5$~ms on integrated GPU).
Execute the interference shader against five regime signatures ($5 \times 0.5 = 2.5$~ms).
Execute the quality metric shader ($\sim\!0.5$~ms).
Read back the regime classification ($\sim\!0.1$~ms).

\item \textbf{Total latency:} $< 4$~ms per classification.
At 2-second windows with 1-second overlap, the system classifies the regime twice per second with $<4$~ms latency.
\end{enumerate}

\subsection{Sleep Stage Detection as Regime Sequence}

Sleep stages map to operational regimes as follows~\cite{sachikonye2026regime,berry2017,steriade1993}:

\begin{center}
\begin{tabular}{lll}
\toprule
Sleep Stage & Regime & $R$ (typical) \\
\midrule
Wake (eyes open) & Turbulent/Aperture & 0.2--0.4 \\
Wake (eyes closed) & Aperture/Cascade & 0.4--0.6 \\
N1 & Cascade & 0.5--0.65 \\
N2 (spindles) & Cascade/Coherent & 0.6--0.85 \\
N3 (slow-wave) & Coherent & 0.8--0.93 \\
REM & Turbulent/Aperture & 0.2--0.5 \\
\bottomrule
\end{tabular}
\end{center}

A full night's hypnogram is the time series $\{k^*(t_j)\}$ of regime classifications.
The GPU pipeline produces this in real time by repeated application of the observation--interference--classification pipeline.

\subsection{Drug-Induced Regime Traversal Monitoring}

Drug action as regime transition~\cite{sachikonye2026aperture,cipriani2018} is monitored by tracking the operator trajectory $\gamma(t)$ in real time.
The structural factor $S(t) = R(t) \exp(-\sigma^2(t)/2\pi^2)$ serves as a scalar summary of therapeutic progress.
The Euler--Lagrange shader (Listing~\ref{lst:euler}) predicts the expected trajectory given drug parameters; deviation between predicted and observed trajectories signals unexpected drug response.

\subsection{Consciousness Window Estimation}

Consciousness in the partition framework is the intersection of decay curves~\cite{sachikonye2026operator}:
\begin{equation}
\boxed{C(t) = P_{\mathrm{decay}}(t) \cap T_{\mathrm{decay}}(t)}
\label{eq:consciousness}
\end{equation}
where $P_{\mathrm{decay}}(t) = P_0 \exp(-t/\tau_P)$ is the pharmacological decay and $T_{\mathrm{decay}}(t) = T_0 \exp(-t/\tau_T)$ is the thermal decay.
The consciousness window is the time interval over which $C(t)$ exceeds a threshold $C_{\min}$:
\begin{equation}
\Delta t_C = \min(\tau_P, \tau_T) \ln\!\frac{\min(P_0, T_0)}{C_{\min}}.
\label{eq:consciousness-window}
\end{equation}

On GPU, $C(t)$ is computed as a per-texel operation: the observation texture at time $t$ is modulated by $P_{\mathrm{decay}}(t)$ and $T_{\mathrm{decay}}(t)$, and the confidence channel (A) encodes whether the consciousness condition is met.

% --- Figure 4 placeholder ---
\begin{figure}[t]
\centering
\fbox{\parbox{0.9\columnwidth}{\centering\vspace{2.5cm}
\textbf{Figure 4:} Sleep stage detection as regime sequence. Top: 8-hour EEG recording (PhysioNet Sleep-EDF~\cite{kemp2000,goldberger2000}) processed through the GPU pipeline. Middle: regime classification $k^*(t)$ color-coded by regime (blue=turbulent, green=aperture, yellow=cascade, orange=coherent, red=phase-locked). Bottom: structural factor $S(t)$ time series. The characteristic NREM--REM cycling is visible as oscillation between coherent and turbulent regimes.
\vspace{2.5cm}}}
\caption{Sleep hypnogram from GPU regime classification.}
\label{fig:sleep}
\end{figure}


% ============================================================================
%                    SECTION 14: HARDWARE REQUIREMENTS
% ============================================================================

\section{Integrated GPU Sufficiency}
\label{sec:hardware}

\begin{theorem}[Memory Bound]
\label{thm:memory}
The complete neural shader-operator pipeline requires at most 25~MB of GPU memory, independent of the recording duration or the number of historical observations.
\end{theorem}

\begin{proof}
We enumerate all GPU-resident allocations:

\begin{enumerate}[leftmargin=*,itemsep=1pt]
\item \textbf{Shader programs:} 5 shaders $\times$ $\sim\!2$~KB compiled $= 10$~KB.
\item \textbf{Observation texture:} $512 \times 512 \times 4 \times 4$ bytes (float32 RGBA) $= 4$~MB.
\item \textbf{Regime signature textures:} $5 \times 512 \times 512 \times 4 \times 1$ byte (RGBA8) $= 5$~MB.
\item \textbf{Interference texture:} $512 \times 512 \times 4 \times 4 = 4$~MB.
\item \textbf{Quality metric texture:} $512 \times 512 \times 4 \times 4 = 4$~MB.
\item \textbf{State texture (Euler--Lagrange):} $2 \times 1 \times 4 \times 4 = 32$~bytes.
\item \textbf{Display texture:} $512 \times 512 \times 4 \times 1 = 1$~MB.
\item \textbf{Uniform buffers:} $\sim\!256$~bytes.
\item \textbf{Vertex buffer:} 0~bytes (fullscreen triangle from vertex ID).
\end{enumerate}

Total: $\sim\!18$~MB.
With double-buffering for the observation and state textures, the total is $\sim\!22$~MB.
Adding overhead for framebuffer objects and driver allocations: $\sim\!25$~MB.

This allocation is \emph{constant}: it does not depend on recording duration, database size, or the number of prior observations.
The Storage Redundancy Theorem~\cite{sachikonye2026gpu} guarantees that any previous observation can be regenerated in $\BigO(1)$ time from the S-entropy coordinates, so storing historical textures is unnecessary.
\end{proof}

\begin{theorem}[Compute Bound]
\label{thm:compute}
The complete pipeline (observation, interference, quality metrics, Euler--Lagrange step) requires $\sim\!0.8 \times 10^9$ floating-point operations per frame, achievable at $>30$~fps on any GPU with $\geq 1$~TFLOPS throughput.
\end{theorem}

\begin{proof}
\textbf{Observation shader:} $512^2 = 262{,}144$ texels, each requiring $\sim\!50$ FLOPs (Listing~\ref{lst:partition}: shell capacity, amplitude, phase, structural factor) $= 13.1 \times 10^6$ FLOPs.

\textbf{Interference shader ($\times 5$):} $262{,}144$ texels $\times$ 20 FLOPs (amplitude multiply, cosine, sum) $\times 5 = 26.2 \times 10^6$ FLOPs.

\textbf{Quality metric shader:} $262{,}144$ texels $\times$ 40 FLOPs (gradients, neighbor reads, coherence) $= 10.5 \times 10^6$ FLOPs.

\textbf{Euler--Lagrange:} 2 texels $\times$ 50 FLOPs $= 100$ FLOPs (negligible).

\textbf{Reduction (visibility, regime):} $\sim\!5 \times 10^6$ FLOPs.

Total: $\sim\!55 \times 10^6$ FLOPs per frame.
At 30~fps: $\sim\!1.65 \times 10^9$ FLOPs/s $= 1.65$~GFLOPS.

Integrated GPU capabilities:
\begin{center}
\begin{tabular}{lr}
\toprule
GPU & TFLOPS \\
\midrule
Intel UHD 620 & 0.44 \\
Intel Iris Xe & 1.69 \\
AMD Radeon Vega 8 & 1.13 \\
Apple M1 (GPU) & 2.60 \\
\bottomrule
\end{tabular}
\end{center}

Even the lowest-end Intel UHD 620 at 0.44~TFLOPS provides $440/1.65 \approx 267\times$ headroom over the 30~fps requirement.
The pipeline is therefore achievable on all integrated GPUs manufactured since approximately 2017.
\end{proof}

\begin{remark}[WebGL 2 Deployment]
The entire pipeline is implementable in WebGL~2~\cite{khronos2023webgl2}, requiring no native application installation.
A standard web browser on any laptop or desktop provides the complete shader-operator trajectory observation system.
\end{remark}


% ============================================================================
%                    SECTION 15: STRUCTURAL FACTOR SHADER
% ============================================================================

\section{Complete Structural Factor Computation}
\label{sec:sf-shader}

For completeness, we present the structural factor as a standalone shader operation that can be composed with any other stage.

\begin{lstlisting}[style=glsl,caption={Structural factor computation shader. Evaluates $S = R\exp(-\sigma^2/2\pi^2)$ as a per-texel operation and outputs it alongside the regime-encoded partition value. The structural factor modulates the amplitude channel of the observation texture.},label=lst:structural]
#version 300 es
precision highp float;

uniform sampler2D u_observationTex;
uniform float u_R;
uniform float u_sigma2;
uniform vec2 u_resolution;

out vec4 fragColor;

const float PI = 3.14159265359;

void main() {
  vec2 uv = gl_FragCoord.xy / u_resolution;
  vec4 obs = texture(u_observationTex, uv);

  // Structural factor
  float S = u_R * exp(-u_sigma2
            / (2.0 * PI * PI));

  // Modulate amplitude by structural factor
  float modAmp = obs.r * S;

  // Regime index as normalized float
  float regimeF = 0.0;
  if (u_R >= 0.95) regimeF = 1.0;
  else if (u_R >= 0.8) regimeF = 0.75;
  else if (u_R >= 0.5) regimeF = 0.5;
  else if (u_R >= 0.3) regimeF = 0.25;

  // Output: R=modulated amplitude,
  //   G=phase, B=regime, A=structural factor
  fragColor = vec4(modAmp, obs.g,
                   regimeF, S);
}
\end{lstlisting}

\begin{proposition}[Structural Factor Monotonicity]
\label{prop:sf-monotonicity}
The structural factor $S$ is monotonically increasing in $R$ and monotonically decreasing in $\sigma^2$.
Along any gradient flow trajectory approaching a coherent or phase-locked fixed point, $S(t)$ is monotonically increasing after an initial transient.
\end{proposition}

\begin{proof}
$\partial S / \partial R = \exp(-\sigma^2/2\pi^2) > 0$ for all $\sigma^2$.
$\partial S / \partial \sigma^2 = -R \exp(-\sigma^2/2\pi^2) / (2\pi^2) < 0$ for $R > 0$.
Along a trajectory approaching a coherent fixed point, $R(t) \to R^* > 0.8$ (increasing) and $\sigma^2(t) \to \sigma^{*2}$ (decreasing), so both partial derivatives push $S$ upward.
\end{proof}


% ============================================================================
%                    SECTION 16: TRAJECTORY VISUALIZATION
% ============================================================================

\section{Trajectory Observation in Real Time}
\label{sec:realtime}

\subsection{Frame-by-Frame Pipeline}

Each frame of the real-time pipeline executes the following sequence:

\begin{enumerate}[leftmargin=*,itemsep=1pt]
\item \textbf{CPU:} Receive 2-second EEG window.
Compute $(\Sk, \St, \Se, R, \sigma^2)$.
Upload as uniforms.

\item \textbf{GPU Pass 1:} Observation shader (Listing~\ref{lst:partition}).
Writes $\Tex_{\mathrm{obs}}$ ($512 \times 512$, RGBA float32).

\item \textbf{GPU Pass 2:} Interference shader ($\times 5$).
Writes $\Tex_{\mathrm{int}}^{(k)}$ for each regime $k$.

\item \textbf{GPU Pass 3:} Quality metric shader.
Reads $\Tex_{\mathrm{obs}}$, writes $\Tex_{\mathrm{qual}}$.

\item \textbf{GPU Pass 4:} Regime classification (Listing~\ref{lst:regime}).
Reads $\Tex_{\mathrm{qual}}$, uniforms.
Writes $\Tex_{\mathrm{class}}$.

\item \textbf{GPU Pass 5:} Euler--Lagrange integration (Listing~\ref{lst:euler}).
Reads $\Tex_{\mathrm{state}}$, writes $\Tex_{\mathrm{state}}'$.

\item \textbf{GPU Pass 6 (optional):} Display shader (HSV colormapping).
\end{enumerate}

\noindent Total: 6 draw calls per frame. At $<1$~ms per draw call on integrated GPU, the pipeline runs at $>60$~fps.

\subsection{Regime Transition Detection}

A regime transition occurs when $k^*(t_j) \neq k^*(t_{j-1})$.
The transition is characterized by:
\begin{enumerate}[leftmargin=*,itemsep=1pt]
\item The departure regime $k^*(t_{j-1})$ and arrival regime $k^*(t_j)$.
\item The structural factor jump $\Delta S = S(t_j) - S(t_{j-1})$.
\item The observation texture difference $\|\Tex_{t_j} - \Tex_{t_{j-1}}\|_{L^2}$.
\item The interference visibility between consecutive textures $\mathcal{V}(\Tex_{t_j}, \Tex_{t_{j-1}})$.
\end{enumerate}
All four quantities are computed on GPU in a single additional draw call.

% --- Figure 5 placeholder ---
\begin{figure}[t]
\centering
\fbox{\parbox{0.9\columnwidth}{\centering\vspace{2.5cm}
\textbf{Figure 5:} Real-time regime transition detection. Top row: four consecutive observation textures spanning a transition from cascade to coherent regime. Bottom: structural factor $S(t)$, interference visibility $\mathcal{V}(t,t-1)$ between consecutive frames, and regime classification $k^*(t)$. The transition is signaled by a sharp drop in $\mathcal{V}$ coinciding with the regime boundary crossing.
\vspace{2.5cm}}}
\caption{Regime transition detection from consecutive observation textures.}
\label{fig:transition}
\end{figure}


% ============================================================================
%                    SECTION 17: DISCUSSION
% ============================================================================

\section{Discussion}
\label{sec:discussion}

\subsection{Relation to Existing Neural Monitoring}

Conventional EEG analysis pipelines compute spectral features (band power ratios, spectral edge frequency) and statistical measures (variance, kurtosis) on the CPU, then apply machine learning classifiers (SVMs, CNNs, LSTMs) trained on labeled data~\cite{delorme2004,nunez2006}.
The free energy principle~\cite{friston2010,friston2006} provides a variational framework but requires substantial computational resources and does not naturally map to GPU observation.
Prior GPGPU approaches~\cite{owens2008,nickolls2008,harris2002} treat the GPU as an accelerator performing the same computations faster, not as an observation apparatus.
The conventional approach has three fundamental limitations: (i)~dependence on labeled training data, (ii)~computational cost scaling with model complexity, and (iii)~opaque classification decisions.

The shader-operator approach eliminates all three.
Classification is performed by physical interference---a deterministic operation with no trainable parameters.
The GPU observables ($\mathcal{P}, \mathcal{C}, \mathcal{V}, \mathcal{R}$) provide objective quality metrics that are interpretable as partition properties.
The computational cost is constant ($\sim\!55$~MFLOP/frame) regardless of the complexity of the neural state.

\subsection{The O=C=P Principle Instantiated}

The neural domain shader instantiates the O=C=P principle (Observation = Computation = Processing)~\cite{sachikonye2026gpu}:

\begin{enumerate}[leftmargin=*,itemsep=1pt]
\item \textbf{Observation.} The partition observation shader performs a categorical measurement of the neural state.
\item \textbf{Computation.} The interference shader computes regime similarity.
The Euler--Lagrange shader computes trajectory dynamics.
\item \textbf{Processing.} The quality metric shader extracts physical observables.
The regime classification shader produces the final output.
\end{enumerate}

All three operations are fragment shader evaluations.
The distinction between observation, computation, and processing collapses: they are the same operation applied to different inputs.

\subsection{Limitations}

The current framework assumes stationary S-entropy within each 2-second analysis window.
For rapidly evolving states (e.g., seizure onset, which can develop over $<1$~second), shorter windows with correspondingly noisier S-entropy estimates may be required.
The explicit Euler integrator in Listing~\ref{lst:euler} may exhibit instability for stiff parameter regimes; a multi-pass RK4 integrator would add three draw calls but improve stability.

The phase-variance relation~\eqref{eq:variance-R} assumes approximately Gaussian phase distributions.
For strongly non-Gaussian distributions (e.g., multimodal phase clustering in the aperture regime), the von Mises concentration parameter $\kappa$ provides a more accurate characterization.


% ============================================================================
%                    SECTION 18: CONCLUSION
% ============================================================================

\section{Conclusion}
\label{sec:conclusion}

We have derived and implemented the complete GPU pipeline for real-time neural partition observation from three axioms.
The main results are:

\begin{enumerate}[leftmargin=*,itemsep=2pt]
\item The neural domain \texttt{computePartitionValue()} function (Listing~\ref{lst:partition}) maps S-entropy coordinates and Kuramoto parameters to a four-channel observation texture encoding amplitude, phase, partition depth, and confidence.

\item The Euler--Lagrange gradient flow equations~\eqref{eq:rdot}--\eqref{eq:sdot} are integrated in a single-pass shader (Listing~\ref{lst:euler}), enabling real-time trajectory computation.

\item Interference-based regime matching (Theorem~\ref{thm:regime-interference}) replaces algorithmic classifiers with physical wave superposition, producing deterministic regime identification from visibility measurements.

\item The structural factor $S = R\exp(-\sigma^2/2\pi^2)$ is a per-texel $\BigO(1)$ operation that provides a single scalar discriminant sharper than $R$ alone (Proposition~\ref{prop:s-discriminant}).

\item Poincar\'{e} backward computation (Theorem~\ref{thm:backward}) requires only negating the time step in the integration shader, enabling treatment planning from specified terminus states.

\item The complete pipeline requires $\sim\!25$~MB GPU memory and $\sim\!1.65$~GFLOPS, achievable on all integrated GPUs since 2017 (Theorems~\ref{thm:memory}--\ref{thm:compute}).
\end{enumerate}

The Trajectory Non-Identity Theorem (Theorem~\ref{thm:nonidentity}) ensures that the GPU-observable trajectory carries irreducible information beyond the endpoint classification.
Two systems at the same regime but with different histories are distinguishable by their memory textures.
This has immediate clinical relevance: different pharmacological paths to the same therapeutic regime are not equivalent, and the shader pipeline makes this non-equivalence directly observable.

The entire framework deploys as a WebGL~2 application in a standard web browser, requiring no specialized hardware, no cloud connectivity, and no proprietary software.
The fragment shader is the observation apparatus.
The texture is the measurement result.
The GPU is the laboratory.


\bibliographystyle{plainnat}
\bibliography{references}

\end{document}
